%! AP Statistics — Unit 4, Lesson 4.7 — Lesson Plan
\documentclass[10pt]{article}
\usepackage{apstats-article}
\usepackage{apstats-boxes}
\usepackage{graphicx}
\graphicspath{{images/}}
\newcommand{\TallMath}[1]{$\displaystyle #1\rule[-1.4em]{0pt}{3.2em}$}
\newcommand{\UnitNumberName}{Unit 4: Probability, Random Variables, and Probability Distributions \quad}
\newcommand{\LessonNumberName}{Lesson 4.7: Introduction to Random Variables and Probability Distributions}

\begin{document}
\begin{center}
    {\Large\bfseries \CourseName: \SchoolYear \\
    {\small\bfseries \UnitNumberName \LessonNumberName}}
\end{center}

\begin{tcolorbox}[colback=sky, colframe=navy, boxrule=0.9pt, arc=2mm,
    left=2mm, right=2mm, top=1.5mm, bottom=1.5mm]
    \textbf{Primary Objective:} Students will represent and interpret probability distributions
    for discrete random variables --- constructing tables and histograms, verifying that all
    probabilities sum to 1, and describing a distribution's shape, center, and spread in context
    (Big Idea: VAR-5).
\end{tcolorbox}

\renewcommand{\arraystretch}{1.6}
\begin{skillbox}[Priority Ideas \& Skills]{goldbox}
\begin{minipage}[t]{0.48\linewidth}
    \textbf{AP Skill 2 --- Data Analysis}
    \begin{itemize}
        \item Represent a probability distribution as a table, graph, or function (Skill 2.B;
              VAR-5.A, VAR-5.A.1--4).
        \item Interpret a probability distribution in context: shape, center, spread (Skill 4.B;
              VAR-5.B, VAR-5.B.1).
        \item Verify that probabilities are valid: each $P(X=x) \ge 0$ and $\sum P(X=x) = 1$.
    \end{itemize}
\end{minipage}
\hfill
\begin{minipage}[t]{0.48\linewidth}
    \textbf{Key Understandings}
    \begin{itemize}
        \item A random variable assigns a number to each outcome of a random process.
        \item A discrete random variable takes a countable number of values; each value carries
              a probability.
        \item A probability distribution completely describes the long-run behavior of a random
              variable.
        \item A cumulative distribution gives $P(X \le x)$ for each value.
        \item Interpreting a distribution means describing shape, center, and spread
              \emph{in context}, not just stating numerical values.
    \end{itemize}
\end{minipage}
\end{skillbox}

\begin{skillbox}[{Vocabulary, Concepts, \& Theorems}]{greenbox}
\begin{minipage}[t]{0.48\linewidth}
    \begin{itemize}
        \item \textbf{Random variable} --- a variable whose value is a numerical outcome of a
              random process
        \item \textbf{Discrete random variable} --- a random variable that takes a countable
              number of values
        \item \textbf{Probability distribution} --- a table, graph, or function giving $P(X=x)$
              for every possible value
    \end{itemize}
\end{minipage}
\hfill
\begin{minipage}[t]{0.48\linewidth}
    \begin{itemize}
        \item \textbf{Probability histogram} --- a bar graph with values of $X$ on the
              horizontal axis and probabilities on the vertical axis
        \item \textbf{Cumulative probability distribution} --- gives $P(X \le x)$ for each $x$
        \item \textbf{Valid probability distribution} --- requires $P(X=x) \ge 0$ for all $x$
              and $\sum P(X=x) = 1$
    \end{itemize}
\end{minipage}
\end{skillbox}

\begin{skillbox}[Activate Prior Knowledge \& Spiral Review]{sky}
    \begin{minipage}[t]{0.58\linewidth}\vspace{0pt}
        \textbf{Prior Knowledge Check (3--4 min)}
        \begin{itemize}
            \item Basic probability rules: $0 \le P(A) \le 1$; $P(A^c) = 1 - P(A)$; addition
                  rule for mutually exclusive events (Unit 4, Lessons 4.1--4.2).
            \item Sample space and outcomes: listing all possibilities for a simple random
                  process.
            \item Prompt: \textit{``If I roll a fair die, list all outcomes and their
                  probabilities. What must those probabilities sum to?''}
        \end{itemize}
        \textbf{Connections Forward}
        \begin{itemize}
            \item Mean/expected value of a distribution $\to$ Lesson 4.8
            \item Combining random variables $\to$ Lesson 4.9
            \item Binomial distribution (a special discrete RV) $\to$ Lesson 4.10
        \end{itemize}
    \end{minipage}
    \hfill
    \begin{minipage}[t]{0.38\linewidth}\vspace{0pt}\centering
        % Authored warm-up; see warmup/main.tex for spiral review problems.
    \end{minipage}
\end{skillbox}

\begin{skillbox}[Hook]{sky}
\begin{minipage}[t]{0.48\linewidth}
    \textbf{Hook (4--5 min)}\\
    Display a table showing the number of pets owned by households in a neighborhood:
    \begin{center}
    \renewcommand{\arraystretch}{1.3}
    \begin{tabular}{@{}ccccc@{}}
    Pets ($x$) & 0 & 1 & 2 & 3 \\
    \hline
    Probability & 0.30 & 0.40 & 0.20 & 0.10 \\
    \end{tabular}
    \end{center}
    Ask: ``If you choose a household at random, what is the probability it owns
    \emph{at most} 1 pet?''\\[4pt]
    Follow-up: ``What must the probabilities sum to, and why?''
\end{minipage}
\hfill
\begin{minipage}[t]{0.48\linewidth}
    \textbf{Lesson Flow (15--18 min)}
    \begin{enumerate}
        \item Define \emph{random variable} and \emph{discrete} vs.\ continuous (brief preview).
        \item Build a probability distribution table from a sample space (number of siblings
              context).
        \item Check the two validity conditions.
        \item Sketch a probability histogram; connect bar heights to probabilities.
        \item Introduce cumulative distribution: $P(X \le x)$ column added to the table.
        \item Interpret the distribution: What does the shape tell us? Where is the center?
              How spread out are the values?
    \end{enumerate}
\end{minipage}
\end{skillbox}

\begin{skillbox}[Explicit Instruction: Building a Probability Distribution]{sky}
\begin{minipage}[t]{0.48\linewidth}
    \textbf{Steps}
    \begin{enumerate}
        \item Identify the random variable $X$ and state its possible values.
        \item Find or assign a probability for each value (from data, symmetry, or a model).
        \item Check: $P(X=x) \ge 0$ for all $x$; $\sum P(X=x) = 1$.
        \item Organize into a two-row table: values in row 1, probabilities in row 2.
        \item Add a cumulative column: $P(X \le x) = \sum_{k \le x} P(X=k)$.
        \item Sketch the probability histogram.
    \end{enumerate}
\end{minipage}
\hfill
\begin{minipage}[t]{0.48\linewidth}
    \textbf{Worked Example}\\
    $X$ = number of cars owned per household (survey of 200 households):
    \begin{center}
    \renewcommand{\arraystretch}{1.3}
    \small
    \begin{tabular}{@{}cccc@{}}
    $x$ & Freq & $P(X=x)$ & $P(X \le x)$ \\
    \hline
    0 & 20  & 0.10 & 0.10 \\
    1 & 80  & 0.40 & 0.50 \\
    2 & 70  & 0.35 & 0.85 \\
    3 & 30  & 0.15 & 1.00 \\
    \end{tabular}
    \end{center}
    Check: $0.10+0.40+0.35+0.15 = 1.00$ \checkmark\\[4pt]
    \textit{Interpretation:} The distribution is roughly symmetric (slight right skew).
    Most households own 1--2 cars.
\end{minipage}
\end{skillbox}

\begin{skillbox}[Active Monitoring]{redbox}
\begin{minipage}[t]{0.48\linewidth}
    \textbf{Circulate and Check}
    \begin{itemize}
        \item Students often confuse frequency with probability --- look for raw counts in the
              $P(X=x)$ column.
        \item Check that cumulative probabilities are non-decreasing and end at exactly 1.
        \item Cold-call: ``How do you know this is a valid probability distribution?''
        \item Cold-call: ``What does $P(X \le 2) = 0.85$ mean in context?''
    \end{itemize}
\end{minipage}
\hfill
\begin{minipage}[t]{0.48\linewidth}
    \textbf{Common Errors}
    \begin{itemize}
        \item Writing probabilities as percentages (e.g., 40 instead of 0.40).
        \item Histogram bars not centered on the values of $X$.
        \item ``Interpreting'' the distribution by only restating the table numbers, not
              describing shape/center/spread in context.
    \end{itemize}
\end{minipage}
\end{skillbox}

\begin{skillbox}[Group Work \& Differentiation]{redbox}
\begin{minipage}[t]{0.48\linewidth}
    \textbf{Activity Summary (10 min)}\\
    Three-tier differentiated activity (see \texttt{activity/}):
    \begin{itemize}
        \item \textbf{Tier R:} Given a frequency table for number of siblings, students
              convert frequencies to probabilities, verify validity, and answer basic
              probability questions.
        \item \textbf{Tier A:} Students add a cumulative column and interpret $P(X \le x)$
              statements in context.
        \item \textbf{Tier E:} Students compare two distributions for number of pets
              owned, describe differences in shape/center/spread, and construct a joint
              table (extension).
    \end{itemize}
\end{minipage}
\hfill
\begin{minipage}[t]{0.48\linewidth}
    \textbf{Differentiation Notes}
    \begin{itemize}
        \item \textit{Support}: Provide a reference card showing the two validity conditions
              and the histogram template with labeled axes.
        \item \textit{ELL}: Vocabulary anchor chart: random variable / distribution /
              probability / cumulative.
        \item \textit{Extension}: Ask students to create their own discrete RV (e.g., number
              of text messages sent per hour) and build a distribution from their own data.
    \end{itemize}
\end{minipage}
\end{skillbox}

\begin{skillbox}[Individual Work \& Assessment]{redbox}
\begin{minipage}[t]{0.48\linewidth}
    \textbf{Exit Ticket (5 min)}\\
    $X$ = number of siblings for a randomly chosen student.
    \begin{center}
    \small
    \renewcommand{\arraystretch}{1.3}
    \begin{tabular}{@{}ccccc@{}}
    $x$ & 0 & 1 & 2 & 3 \\
    \hline
    $P(X=x)$ & 0.15 & 0.45 & 0.30 & 0.10 \\
    \end{tabular}
    \end{center}
    \begin{enumerate}
        \item Is this a valid probability distribution? Explain.
        \item Find $P(X \ge 2)$.
        \item Describe the shape of this distribution in context.
    \end{enumerate}
\end{minipage}
\hfill
\begin{minipage}[t]{0.48\linewidth}
    \textbf{AP-Style Check}\\
    \textit{Which of the following could be a valid probability distribution for
    a discrete random variable?}
    \begin{enumerate}[label=(\Alph*)]
        \item $P(0)=0.30,\ P(1)=0.40,\ P(2)=0.40$
        \item $P(0)=0.25,\ P(1)=0.50,\ P(2)=0.25$
        \item $P(0)=-0.10,\ P(1)=0.60,\ P(2)=0.50$
        \item $P(0)=0.10,\ P(1)=0.40,\ P(2)=0.60$
    \end{enumerate}
    Answer: (B) --- only option with all non-negative probabilities summing to 1.
\end{minipage}
\end{skillbox}

\begin{skillbox}[Reinforcement \& Extension]{goldbox}
\begin{minipage}[t]{0.48\linewidth}
    \textbf{Homework / Practice}
    \begin{itemize}
        \item Six problems: build distributions from frequency tables, verify validity,
              compute $P(X \le x)$ and $P(X > x)$, and interpret in context.
        \item Emphasis on full sentence interpretations (AP exam rewards context).
    \end{itemize}
\end{minipage}
\hfill
\begin{minipage}[t]{0.48\linewidth}
    \textbf{Extension / Preview}
    \begin{itemize}
        \item \textit{Extension}: Given a probability distribution, find the missing
              probability and explain why it must take that value.
        \item \textit{Preview}: Lesson 4.8 --- once we have a distribution, how do we
              describe its average (expected value) and typical spread numerically?
    \end{itemize}
\end{minipage}
\end{skillbox}

\end{document}
